\section{Related work}\label{sec:related}
Our results draw mainly on two bodies of work. Research on indicator
quadratics with sparse matrices supplies the model, the messages of
Section~\ref{sec:model}, and the structures under which exact or
approximate methods are known to be efficient. Smoothed analysis of
discrete optimization supplies the counting and scaling arguments behind
Theorem~\ref{thm:main} and Corollary~\ref{cor:magnitude}.
Sections~\ref{sec:hardness}--\ref{sec:discussion} give further detail
next to individual results.

\subsection{Exact algorithms and polynomial cases}
For tridiagonal matrices, whose support graph is a path, Liu, Fattahi,
G\'omez, and K\"u\c{c}\"ukyavuz solve the problem in quadratic time by
dynamic programming and give a compact extended formulation of the
closed convex hull of the epigraph \cite{LiuPath}. Bhathena, Fattahi,
G\'omez, and K\"u\c{c}\"ukyavuz \cite{BhathenaTrees} solve the problem
on arbitrary trees, also in quadratic arithmetic time and memory. Their
dynamic program passes parametric costs and their conjugates, and our
messages play the role of their parametric costs.
Theorem~\ref{thm:main} does not improve this tree-specific guarantee.

For graphs of larger treewidth, the same authors build exact parametric
dynamic programs over a tree decomposition \cite{BhathenaStructured}.
Their local parametric costs, defined for a whole bag, are minima of one
quadratic per support, like our messages \cite[Section~2.2 and
Lemma~1]{BhathenaStructured}. They call these quadratics quadratic
pieces; we call them branches. Their costs also include the quadratic
and linear terms of the bag variables, so their quadratics are strongly
convex rather than concave; these terms are the same for every internal
support and do not change which supports are optimal. Their dynamic
program computes these costs bottom-up, as outlined in
Section~\ref{sec:prelim-messages}, and after each step discards the
branches that a pairwise test shows are not optimal anywhere on a box of
boundary values that contains the optimal solutions
\cite[Definition~3 and Section~5]{BhathenaStructured}.

The running-time bound of Bhathena et al.\ depends on the treewidth, the
conditioning of $Q$, the volume growth of the support graph, a
deterministic margin $(k,\eta)$, and the radius $U$ of the box on which
they prune, which they take proportional to $\|c\|_\infty$
\cite[Lemma~3, Definition~5, Lemma~9, and
Theorem~1]{BhathenaStructured}. Volume growth measures how fast the
number of vertices within a given graph distance of a bag grows with
that distance, and the margin bounds by $k$ the number of
$\eta$-optimal local supports. When all of these are fixed, the bound is
linear in $n$. With linear volume growth it is polynomial in
$U^2/\eta$, so, like Theorem~\ref{thm:main}, it grows with the size of
the linear coefficients \cite[Theorem~1 and
Corollary~1]{BhathenaStructured}.

Their pruning compares only supports that agree on the vertices within
a certain graph distance of a bag and keeps at most $k$ of them for
each such pattern; volume growth bounds the number of these vertices,
and hence of the patterns \cite[Lemmas~8 and~9]{BhathenaStructured}.
Theorem~\ref{thm:main} needs neither the margin nor volume growth: under
penalty perturbation, Lemma~\ref{lem:count} bounds the expected number
of near-optimal supports of a whole subproblem at each net point. On the
other hand, Theorem~\ref{thm:main} concerns the perturbed problem, and
its bound is a polynomial of high degree. Neither result implies the
other. Section~\ref{sec:discussion} estimates the size of our bound,
states their margin condition precisely, and explains how the
perturbation affects it.

Other polynomial cases rest on special matrix structure. Problems with a
Stieltjes matrix, whose off-diagonal entries are nonpositive, and with
linear coefficients of one sign can be solved in polynomial time
\cite{AtamturkGomezM}, as recalled in \cite{LiuStieltjes}. Lee,
G\'omez, and Atamt\"urk study multi-period problems in which the states
follow exact affine dynamics, the quadratic part of the cost involves
only the states, and only the controls carry indicators. Projecting out
the states then leaves a positive-definite matrix with a (block)
tridiagonal inverse \cite[Propositions~1 and~2]{LeeMultiperiod}. They
use this structure to describe the closed convex hull of the epigraph
\cite[Theorems~1 and~2]{LeeMultiperiod} and, when there are no side
constraints, to solve the problem as a shortest-path problem on a
directed acyclic graph \cite[Proposition~9 and
Corollary~3]{LeeMultiperiod}. G\'omez, Han, and Lozano observe that the
tridiagonal, tree, and Stieltjes structures behind earlier polynomial
cases disappear for bandwidth greater than one
\cite[Section~1.1]{GomezBanded}. Section~\ref{sec:hardness} shows that
exact optimization is NP-hard already at bandwidth two, even for
matrices arbitrarily close to the identity, and that its hard instances
cannot be made Stieltjes by changing the signs of variables.

\subsection{Approximation, decision diagrams, and convex hulls}
Approximation schemes also exploit bounded treewidth. The linear
programs of Bienstock and Mu\~noz \cite{BienstockMunoz}
(Section~\ref{sec:prelim-treedec}), of polynomial size for each fixed
tolerance, replace each continuous variable by a truncated binary
expansion. Bienstock and Chen \cite[Theorem~1.4]{BienstockChen} treat
convex quadratic problems in which binary indicators control disjoint
blocks of variables, when the graph whose vertices are the constraints,
with two constraints adjacent if they involve a common block, has
bounded treewidth. Their algorithm returns a solution that costs at
most the optimal value, satisfies the constraints on the indicators
alone exactly, and violates the other constraints by a bounded amount.
Like the linear programs of Bienstock and Mu\~noz, the grid oracle of
Lemma~\ref{lem:gridoracle} rounds every coordinate to a grid; it then
solves the rounded problem by dynamic programming on the tree
decomposition.
Theorem~\ref{thm:dictionary} accepts any restricted additive oracle in
its place.

Decision diagrams take a different route. A decision diagram is a
layered directed acyclic graph whose layers correspond to the indicator
variables. Each path from the root to a terminal node encodes one
support; in an exact diagram, the length of the path is the cost of that
support, so a shortest path solves the problem
\cite[Section~2.1]{GomezBanded}. The nodes are the states of a dynamic
program that fixes the indicators one at a time and stores what the rest
of the problem needs to know about the partial support. A message, in
contrast, conditions on continuous boundary values and minimizes over
the internal supports. G\'omez, Han, and Lozano \cite{GomezBanded} build
decision diagrams for banded matrices, derive from them extended
formulations of the closed convex hull and, from truncated diagrams, an
additive approximation scheme, and use the diagrams to re-solve online
instances quickly. Choi et al.\ \cite{ChoiDiagrams} develop decision
diagrams that both solve indicator quadratic problems and give ideal
conic extended formulations. Their exact diagrams have polynomial size
for fixed rank and for trees and inverse trees with a fixed number of
leaves, with $O(n^{k+1})$ nodes for $k$ leaves
\cite[Section~7]{ChoiDiagrams}. For each fixed accuracy, their
approximate diagrams have size linear in $n$ under polynomial volume
growth and a polynomial bound on how many already decided variables lie
within a given graph distance of the undecided ones
\cite[Definition~15 and Theorem~4]{ChoiDiagrams}.
Section~\ref{sec:discussion} explains how such approximate diagrams can
serve as restricted additive oracles.

The diagrams of \cite{GomezBanded} and \cite{ChoiDiagrams} are built
from the matrix alone, and only their arc lengths depend on the linear
coefficients and penalties \cite[Remark~1]{GomezBanded},
\cite[Definition~7 and Proposition~6]{ChoiDiagrams}. A representation of
this kind can be built once and reused for every linear cost and penalty
vector. A message can be reused when data outside its subproblem change
(Section~\ref{sec:prelim-chain}), but not for every cost vector, because
it depends on the linear coefficients and perturbed penalties of its
internal variables.

Representations built from $Q$ alone can be exponentially large even
where optimization is easy. Choi et al.\ observe this for their own
diagrams: on stars, trees in which one center vertex is adjacent to all
$n-1$ other vertices, so that the number of leaves grows linearly with
$n$, their representation becomes exponential in size
\cite[Section~7.2]{ChoiDiagrams}, and Remark~\ref{rem:dd-star} checks
this directly for their state definition and every variable order.
Appendix~\ref{sec:representation} shows the same for the
inverse-principal polytope $P_Q$, built from the inverses of the
principal submatrices of $Q$, through which Wei et al.\
\cite[Section~3.1 and Theorem~1]{WeiHull} describe the closed convex
hull of the epigraph of $x^TQx$ under the indicator constraints. For a
sparse, well-conditioned star, every exact linear, second-order-cone, or
fixed-order semidefinite lift of $P_Q$ has exponential size
(Corollary~\ref{cor:conic}). This bound concerns $P_Q$ together with its
inverse entries; it gives no lower bound for formulations of the
epigraph hull that project them out.

\subsection{Smoothed analysis and enumeration}
Theorem~\ref{thm:main} is a result of smoothed analysis in the sense of
Section~\ref{sec:intro}, with arbitrary data and random perturbations of
the penalties only. For binary problems with perturbed linear
objectives, two principles used here are already established: exact
optimization in expected polynomial time under perturbation, and the
conversion of an approximation scheme into such exact optimization. A
binary optimization problem has \emph{polynomial smoothed complexity} if
some algorithm solves it with running time $T$ satisfying
$\E[T^\alpha]\le\beta\phi\ell$ for constants $\alpha,\beta>0$, on every
instance of input length $\ell$ whose perturbed coefficients have
densities at most $\phi\ge1$ \cite[conference version,
Section~1.3]{BeierVocking}; this is weaker than expected polynomial
running time. Beier and V\"ocking \cite[conference version,
Theorem~3]{BeierVocking} showed that such a problem has polynomial
smoothed complexity if and only if it has a randomized pseudopolynomial
algorithm, and R\"oglin and Teng \cite[Section~6.2]{RoglinTeng}
strengthened this to expected polynomial running time. Dughmi and
Roughgarden \cite[Proposition~2.4]{DughmiRoughgarden} explicitly convert
an approximation scheme into exact smoothed optimization for binary
linear objectives.

Lemma~\ref{lem:count} combines the sandwich inequality, which bounds the
size of a family of supports by the number of coordinate sets on which
it shows all binary patterns \cite[Theorem~5]{KozmaMoran}, with the
conditioning argument of Beier and V\"ocking's generalized isolating
lemma \cite[conference version, Lemma~5]{BeierVocking}
(Section~\ref{sec:enumeration}). R\"oglin and Teng
\cite[Lemma~6.1]{RoglinTeng} bound, for a linear objective whose
coefficients have bounded densities, the
probability that the $u$th best solution is within $\varepsilon$ of the
optimum, that is, the probability that at least $u$ solutions lie within
$\varepsilon$ of the optimum. For linear costs this is a tail bound for
the count in Lemma~\ref{lem:count}, proved by a related conditioning
argument. Lemma~\ref{lem:count} bounds the expectation instead, allows
an arbitrary deterministic cost for each support, and allows
perturbation values of positive probability (atoms), which the finite
perturbation grid requires.

Smoothed analysis has also been applied to outputs that, like our
dictionaries, describe many objectives at once. Beier, R\"oglin,
R\"osner, and V\"ocking \cite[Theorem~1]{BeierPareto} bound expected
bicriteria Pareto counts with one arbitrary objective and one linear
objective whose independent coefficients have bounded densities and
finite first moments. Moments of smoothed Pareto counts are also known
\cite{RoglinTeng,BrunschRoglin}, and smoothed parametric enumeration has
been done for shortest paths \cite[Section~6]{ApplegatePaths}. These
results use different perturbation models and outputs.
Theorem~\ref{thm:dictionary} allows arbitrary deterministic Lipschitz
branches and adds $\xi^Tz$ to the branch of each support $z$, with one
random $\xi$ shared by all supports and parameter values. The
ingredients above are classical. To the best of our knowledge, what is
new is the two-sided enumeration guarantee for a restricted additive
oracle, which adapts Lawler's partition method \cite{Lawler}
(Lemma~\ref{lem:enumeration}), its use with a parameter net to
obtain complete dictionaries on a whole parameter domain
(Theorem~\ref{thm:dictionary}), and the specialization of this
construction to indicator quadratics (Theorem~\ref{thm:main}).

\subsection{Lower bounds}
Chains that encode SUBSET SUM through a continuous linear recursion
already occur in conditional linear Gaussian inference
\cite[Section~3]{LernerParr}. Theorem~\ref{thm:hardness} adapts this
mechanism to one positive-definite quadratic with uncoupled indicators
and positive penalties, at bandwidth two and near-identity conditioning,
in three numerical regimes. Corollary~\ref{cor:magnitude} carries the
hardness over to the perturbed problem by a scaling argument of Beier
and V\"ocking \cite[conference version, Theorem~3]{BeierVocking}.

Hardness of sparse approximation is also known in general. \c{C}ivril
\cite[Theorem~1.1]{Civril} proves inapproximability for a
cardinality-constrained squared-projection objective without graph,
positive-definiteness, or spectral restrictions. In the other direction,
Das and Kempe \cite[Theorem~4.1]{DasKempe} give a fully polynomial-time
approximation scheme for cardinality-constrained subset selection when
the covariance graph of the predictors alone, which corresponds to our
matrix graph, has bounded bandwidth. Their guarantee is a relative
approximation of explained variance, so it does not conflict with exact
hardness of the penalized problem. For fixed bandwidth and eigenvalue
bounds, the additive approximation scheme of G\'omez, Han, and Lozano
runs in time polynomial in the dimension, a bound on the linear
coefficients, and the inverse accuracy $1/\varepsilon$
\cite[Theorem~1]{GomezBanded}. Case~(1) of Theorem~\ref{thm:hardness},
with a fixed matrix family, exponentially large coefficients, and a
constant gap between the optimal values of YES and NO instances, shows
that its dependence on the coefficient bound cannot be replaced by
dependence on encoding length unless $\mathrm{P}=\mathrm{NP}$.
Case~(2), with $|c_i|\le2$ and a possibly exponentially small gap,
shows the same for its dependence on $1/\varepsilon$.

Value functions with exponentially many pieces are known in other
settings. For parametric linear programs, the optimal value can have
exponentially many linear pieces \cite[Corollary~1]{Murty}, and
exponential growth of piecewise quadratic messages is known for broader
nonconvex models \cite[Section~5]{KuricTree}.
Theorem~\ref{thm:fixed-data} realizes such growth in the indicator model
with fixed local data and near-identity conditioning. Power-law rates
like that of Theorem~\ref{thm:pruning} are known for approximation by
quadratic bases with a common Hessian \cite[Section~III,
Theorem~3.3]{GaubertPruning}, and pruning under uniform error has a
covering interpretation \cite[Section~VI]{GaubertPruning}.
Theorem~\ref{thm:pruning} gives matching upper and lower exponents for
branches whose Hessians differ between supports.
